% 本模板改编自中国科学院大学本科生公共必修课程《基础物理实验》Word模板格式编写

\documentclass[11pt]{article}

\usepackage[a4paper]{geometry}
\geometry{left=2.0cm,right=2.0cm,top=2.5cm,bottom=2.5cm}

\usepackage{ctex} % 支持中文的LaTeX宏包
\usepackage{amsmath,amsfonts,graphicx,subfigure,amssymb,bm,amsthm,mathrsfs,mathtools,breqn} % 数学公式和符号的宏包集合
\usepackage{algorithm,algorithmicx} % 算法和伪代码的宏包
\usepackage[noend]{algpseudocode} % 算法和伪代码的宏包
\usepackage{fancyhdr} % 自定义页眉页脚的宏包
\usepackage[framemethod=TikZ]{mdframed} % 创建带边框的框架的宏包
\usepackage{fontspec} % 字体设置的宏包
\usepackage{adjustbox} % 调整盒子大小的宏包
\usepackage{fontsize} % 设置字体大小的宏包
\usepackage{tikz,xcolor} % 绘制图形和使用颜色的宏包
\usepackage{multicol} % 多栏排版的宏包
\usepackage{multirow} % 表格中合并单元格的宏包
\usepackage{pdfpages} % 插入PDF文件的宏包
\RequirePackage{listings} % 在文档中插入源代码的宏包
\RequirePackage{xcolor} % 定义和使用颜色的宏包
\usepackage{wrapfig} % 文字绕排图片的宏包
\usepackage{bigstrut,multirow,rotating} % 支持在表格中使用特殊命令的宏包
\usepackage{booktabs} % 创建美观的表格的宏包
\usepackage{circuitikz} % 绘制电路图的宏包
\usepackage{float}
\usepackage{siunitx}

\definecolor{dkgreen}{rgb}{0,0.6,0}
\definecolor{gray}{rgb}{0.5,0.5,0.5}
\definecolor{mauve}{rgb}{0.58,0,0.82}
\lstset{
  frame=tb,
  aboveskip=3mm,
  belowskip=3mm,
  showstringspaces=false,
  columns=flexible,
  framerule=1pt,
  rulecolor=\color{gray!35},
  backgroundcolor=\color{gray!5},
  basicstyle={\small\ttfamily},
  numbers=none,
  numberstyle=\tiny\color{gray},
  keywordstyle=\color{blue},
  commentstyle=\color{dkgreen},
  stringstyle=\color{mauve},
  breaklines=true,
  breakatwhitespace=true,
  tabsize=3,
}

% 轻松引用, 可以用\cref{}指令直接引用, 自动加前缀.
% 例: 图片label为fig:1
% \cref{fig:1} => Figure.1
% \ref{fig:1}  => 1
\usepackage[capitalize]{cleveref}
% \crefname{section}{Sec.}{Secs.}
\Crefname{section}{Section}{Sections}
\Crefname{table}{Table}{Tables}
\crefname{table}{Table.}{Tabs.}

% \setmainfont{Palatino Linotype.ttf}
% \setCJKmainfont{SimHei.ttf}
% \punctstyle{kaiming}
% 偏好的几个字体, 可以根据需要自行加入字体ttf文件并调用，如果编译报错找不到字体，请注释掉上述三行使用默认字体

\renewcommand{\emph}[1]{\begin{kaishu}#1\end{kaishu}}

%改这里可以修改实验报告表头的信息
\newcommand{\experiName}{基于ResNet18迁移学习的手势姿态识别}
\newcommand{\supervisor}{Jun-Xiong Ji}
\newcommand{\name}{史峰源}
\newcommand{\studentNum}{2412526}
\newcommand{\major}{智能科学与技术}
\newcommand{\room}{xxx}
\newcommand{\others}{$\square$}
%% 如果是调课、补课, 改为: $\square$\hspace{-1em}$\surd$
%% 否则, 请用: $\square$
%%%%%%%%%%%%%%%%%%%%%%%%%%%

\begin{document}
\begin{center}
    \LARGE \bf \, 机\, 器\, 视\, 觉\, 实\, 验\, 报\, 告
\end{center}

\begin{center}
    \noindent
    \emph{专业}\underline{\makebox[8em][c]{\major}}
    \emph{姓名}\underline{\makebox[6em][c]{\name}}
    \emph{学号}\underline{\makebox[6em][c]{\studentNum}}\\
\vspace{0.2em}
    {\noindent}
    \rule[8pt]{17cm}{0.2em}
\end{center}

\begin{center}
    \Large \bf \experiName
\end{center}

\section{实验目的}
\begin{enumerate}
    \item 完成六类手势姿态图像分类任务，识别类别包括A、B、C、Five、Point、V。
    \item 使用ResNet18预训练模型进行迁移学习，构建适用于小规模手势数据集的分类模型。
    \item 实现分层数据划分，使训练集和验证集在各类别上保持一致比例。
    \item 统一验证与推理阶段的确定性预处理流程，保证实验结果具有可复现性。
    \item 通过数据分布图、样例图和训练指标对实验过程进行辅助分析。
\end{enumerate}

\section{实验环境与数据集}

\subsection{实验环境}
本实验使用Python与PyTorch完成模型训练和推理。主要依赖包括torch、torchvision、pillow、tqdm、scikit-learn和matplotlib。训练时优先使用CUDA设备；本机实验环境可识别NVIDIA GeForce RTX 4050 Laptop GPU，因此训练过程使用GPU加速。

\begin{table}[H]
    \centering
    \caption{主要实验配置}
    \begin{tabular}{cc}
        \toprule
        项目 & 配置 \\
        \midrule
        深度学习框架 & PyTorch \\
        基础模型 & ResNet18 \\
        预训练权重 & ImageNet预训练权重 \\
        输入尺寸 & $224 \times 224$ \\
        批大小 & 32 \\
        训练轮数 & 30 \\
        学习率 & $1\times10^{-4}$ \\
        优化器 & AdamW \\
        学习率调度 & CosineAnnealingLR \\
        损失函数 & CrossEntropyLoss，label smoothing为0.1 \\
        \bottomrule
    \end{tabular}
\end{table}

\subsection{数据集组成与划分}
数据集共包含六个类别：A、B、C、Five、Point、V。原始数据按照ImageFolder格式组织，即每个类别对应一个独立文件夹。实验采用分层划分策略，将每一类图像按约80\%和20\%分别划入训练集和验证集，从而使验证集能够保持与原始数据一致的类别比例。

\begin{table}[H]
    \centering
    \caption{各类别图像数量统计}
    \begin{tabular}{cccc}
        \toprule
        类别 & 原始数据 & 训练集 & 验证集 \\
        \midrule
        A & 1327 & 1062 & 265 \\
        B & 508 & 406 & 102 \\
        C & 604 & 483 & 121 \\
        Five & 707 & 566 & 141 \\
        Point & 1394 & 1115 & 279 \\
        V & 500 & 400 & 100 \\
        \midrule
        合计 & 5040 & 4032 & 1008 \\
        \bottomrule
    \end{tabular}
\end{table}

\begin{figure}[H]
     \centering
     \includegraphics[width=0.88\textwidth]{outputs/visualizations/class_distribution.png}
    \caption{原始数据、训练集和验证集的类别数量分布}
    \label{fig:class-distribution}
 \end{figure}

\section{实验原理与方法}

\subsection{迁移学习模型}
ResNet通过残差连接缓解深层神经网络训练中的梯度退化问题，其基本思想是学习残差映射：
\begin{equation}
    \mathcal{H}(x)=\mathcal{F}(x)+x
\end{equation}
其中$x$为输入特征，$\mathcal{F}(x)$为残差分支学习到的变换。相较于从零开始训练卷积神经网络，使用ImageNet预训练的ResNet18能够利用已有的通用视觉特征，如边缘、纹理和局部形状，从而提高小规模图像分类任务的训练效率与稳定性。

本实验保留ResNet18的卷积特征提取部分，并将原始全连接分类层替换为适用于六分类任务的新分类头。分类头包含Dropout、线性层、ReLU激活和最终输出层，用于降低过拟合风险并输出六个类别的logits。

\subsection{数据预处理}
训练阶段使用随机增强以提升模型对尺度、平移、旋转和光照变化的适应能力。具体包括随机裁剪、水平翻转、随机旋转、颜色扰动和仿射变换。验证阶段不再使用随机操作，而是采用确定性的\texttt{Resize(256)}和\texttt{CenterCrop(224)}，保证同一张图像每次评估时输入一致。

训练阶段预处理流程如下：
\begin{lstlisting}[language=Python]
transforms.Resize(256)
transforms.RandomResizedCrop(224, scale=(0.75, 1.0))
transforms.RandomHorizontalFlip(p=0.5)
transforms.RandomRotation(15)
transforms.ColorJitter(
    brightness=0.25,
    contrast=0.25,
    saturation=0.25,
    hue=0.05
)
transforms.RandomAffine(
    degrees=0,
    translate=(0.08, 0.08),
    scale=(0.9, 1.1)
)
transforms.ToTensor()
transforms.Normalize(mean=[0.485, 0.456, 0.406],
                     std=[0.229, 0.224, 0.225])
\end{lstlisting}

验证阶段预处理流程如下：
\begin{lstlisting}[language=Python]
transforms.Resize(256)
transforms.CenterCrop(224)
transforms.ToTensor()
transforms.Normalize(mean=[0.485, 0.456, 0.406],
                     std=[0.229, 0.224, 0.225])
\end{lstlisting}

\subsection{训练策略}
模型训练使用交叉熵损失函数，并加入label smoothing以缓解模型过度自信的问题。优化器选用AdamW，学习率为$1\times10^{-4}$，权重衰减为$1\times10^{-4}$。学习率调度器采用余弦退火策略，在30个epoch内逐步调整学习率。

训练过程中记录每轮训练损失、训练准确率、验证损失和验证准确率。当验证准确率超过历史最佳值时，保存当前模型权重及类别映射信息。保存的checkpoint包含\texttt{model\_state\_dict}、\texttt{class\_to\_idx}、\texttt{classes}和\texttt{best\_acc}，从而避免推理阶段出现类别顺序不一致的问题。

\section{程序代码}

\textit{以下给出本实验中与数据划分、模型构建、训练保存和推理标签映射相关的核心代码片段。完整代码见项目目录中的Python脚本。}

\subsection{数据分层划分}
数据划分脚本按照类别文件夹分别打乱并拆分图像，保证每个类别在训练集和验证集中都保持近似80/20的比例。
\begin{lstlisting}[language=Python]
def split_class_files(files, val_ratio, rng):
    files = list(files)
    rng.shuffle(files)

    if len(files) <= 1:
        return files, []

    val_count = round(len(files) * val_ratio)
    val_count = min(max(1, val_count), len(files) - 1)

    val_files = files[:val_count]
    train_files = files[val_count:]

    return train_files, val_files
\end{lstlisting}

\subsection{模型构建}
模型基于torchvision提供的ResNet18，并替换最后的分类层以适配六分类任务。
\begin{lstlisting}[language=Python]
def build_resnet18(num_classes=6, pretrained=True):
    if pretrained:
        weights = models.ResNet18_Weights.DEFAULT
    else:
        weights = None

    model = models.resnet18(weights=weights)
    in_features = model.fc.in_features
    model.fc = nn.Sequential(
        nn.Dropout(0.4),
        nn.Linear(in_features, 256),
        nn.ReLU(inplace=True),
        nn.Dropout(0.3),
        nn.Linear(256, num_classes)
    )
    return model
\end{lstlisting}

\subsection{训练与checkpoint保存}
训练脚本使用验证准确率选择最优模型，并将类别映射与模型权重一起保存。
\begin{lstlisting}[language=Python]
if val_acc > best_acc:
    best_acc = val_acc
    best_model_wts = copy.deepcopy(model.state_dict())

    torch.save({
        "model_state_dict": best_model_wts,
        "class_to_idx": train_dataset.class_to_idx,
        "classes": train_dataset.classes,
        "best_acc": best_acc
    }, "checkpoints/best_model.pth")
\end{lstlisting}

\subsection{推理阶段类别映射}
推理阶段不再硬编码类别顺序，而是从checkpoint读取训练时保存的类别映射。
\begin{lstlisting}[language=Python]
def get_idx_to_class(checkpoint):
    if "class_to_idx" in checkpoint:
        idx_to_class = {
            int(index): class_name
            for class_name, index in checkpoint["class_to_idx"].items()
        }
    elif "classes" in checkpoint:
        idx_to_class = {
            index: class_name
            for index, class_name in enumerate(checkpoint["classes"])
        }
    else:
        raise ValueError("Checkpoint must include class metadata.")

    return idx_to_class
\end{lstlisting}

\section{实验结果与分析}

\subsection{数据可视化结果}
\Cref{fig:class-distribution}展示了各类别图像数量分布。可以看出，A和Point类别样本数较多，而B和V类别样本数较少，数据集存在一定类别不均衡。通过分层划分，训练集和验证集保留了原始数据的类别比例，使验证结果更能反映同源数据分布下的模型表现。

\begin{figure}[H]
     \centering
     \includegraphics[width=0.86\textwidth]{outputs/visualizations/raw_samples.png}
    \caption{原始数据集中各类别样例}
    \label{fig:raw-samples}
 \end{figure}

\begin{figure}[H]
     \centering
     \includegraphics[width=0.86\textwidth]{outputs/visualizations/train_samples.png}
    \caption{训练集中各类别样例}
    \label{fig:train-samples}
 \end{figure}

\begin{figure}[H]
     \centering
     \includegraphics[width=0.86\textwidth]{outputs/visualizations/val_samples.png}
    \caption{验证集中各类别样例}
    \label{fig:val-samples}
 \end{figure}

\subsection{训练结果}
模型训练共进行30个epoch。训练过程中，模型在验证集上取得的最佳准确率为：
\begin{equation}
    \text{Best Val Acc}=1.0000
\end{equation}
即在当前分层验证集上达到100\%的分类准确率。最佳模型保存为\texttt{checkpoints/best\_model.pth}。从实验结果看，迁移学习能够在较短训练轮数内充分利用预训练特征，并快速适应当前六类手势分类任务。

为进一步观察最终模型在课程验证集上的分类情况，使用保存的最佳checkpoint对\texttt{data/val}中的1008张验证图像重新进行推理，并绘制混淆矩阵和各类别准确率图。结果显示，六个类别的预测均集中在混淆矩阵对角线上，未出现跨类别误判；各类别准确率均为1.0000。

\begin{figure}[H]
     \centering
     \includegraphics[width=0.76\textwidth]{outputs/validation_evaluation/validation_confusion_matrix.png}
    \caption{最终模型在验证集上的混淆矩阵}
    \label{fig:validation-confusion-matrix}
 \end{figure}

\begin{figure}[H]
     \centering
     \includegraphics[width=0.86\textwidth]{outputs/validation_evaluation/validation_per_class_accuracy.png}
    \caption{最终模型在验证集上的各类别准确率}
    \label{fig:validation-per-class-accuracy}
 \end{figure}

需要注意的是，验证集由同一原始数据集分层划分得到，因此该指标主要反映模型在同源数据分布下的识别能力。为了保证实验目标明确，本报告仅分析课程数据集划分得到的训练与验证结果。

\subsection{结果分析}
从数据分布和样例图可以观察到，不同类别在手形、手指数量和姿态方向上具有较明显差异，这有利于卷积神经网络进行特征学习。训练阶段引入随机裁剪、旋转、颜色扰动和仿射变换，可以减少模型对固定位置和固定光照的依赖。验证阶段使用确定性中心裁剪，使每次验证输入一致，保证实验指标可复现。

此外，checkpoint中保存类别映射是本实验中的重要工程改进。若仅在推理代码中手动写入类别列表，一旦训练数据目录顺序或类别名称发生变化，模型输出下标与类别名可能不一致。将\texttt{class\_to\_idx}与模型权重一起保存，可以保证推理阶段与训练阶段使用完全相同的标签语义。

\section{实验结论}
本实验完成了基于ResNet18迁移学习的六类手势姿态识别模型。通过分层划分数据集、统一验证与推理预处理、保存类别映射和可视化数据分布，项目形成了较完整的训练与分析流程。实验结果表明，模型在课程数据集划分得到的验证集上达到100\%准确率，说明迁移学习方法能够有效适应该手势分类任务。

后续若继续改进，可在不改变当前最佳模型的前提下，增加更细粒度的类别混淆分析、训练日志曲线记录和多次随机种子实验，以进一步评估模型稳定性。

\end{document}
